% 第2章问题分析中的问题一片段
\section{问题分析}

\subsection{问题一分析}

问题一给定无人机的飞行方向、飞行速度、投放时刻和起爆延迟，要求计算烟幕对来袭导弹视线的
有效遮蔽时间。其核心并非只判断烟幕球心是否靠近一条空间直线，而是从导弹这一移动观察点出发，
判断真目标在视野中的完整投影能否被烟幕形成的遮挡区域覆盖。

为保留导弹、烟幕和目标之间的远近关系，本文在每一时刻以导弹位置为视点，以导弹指向真目标
几何中心的视线为法向量建立观测平面。随后构造随时间变化的正交视图矩阵，将目标圆柱和烟幕球
转换到以导弹为原点的相机坐标系，再通过深度归一化完成三维到二维的透视投影。普通正投影会忽略
物体距导弹的深度差，从而不能正确反映烟幕和目标的视角大小，因此不作为最终判据。

在二维观测平面上，真目标形成投影集合，烟幕球则由导弹到球面的切锥形成一个二次曲线遮挡区域。
若目标投影完全位于该遮挡区域内，并且每条对应射线均先经过烟幕、后到达目标，则判定该时刻形成
有效遮蔽。由此，问题一最终转化为求上述投影包含关系在烟幕有效期内成立的时间区间及其长度。
